% SPDX-License-Identifier: LPPL-1.3c
% Copyright (C) 2026 Christophe Jorssen
% Author and Current Maintainer: Christophe Jorssen <christophe.jorssen@gmail.com>
% LPPL maintenance status: maintained
% This file is part of luafemm. See LICENSE and LICENSES.md for its terms.
\documentclass[border=3mm]{standalone}
\usepackage{tikz}
\usetikzlibrary{femm}
\begin{document}
\tikzset{paired edge/.style={purple,dashed,thick,
  femm/frontier={name=#1,type=periodic}},
  femm/materials/recoil/.style={mu r=1.05,coercivity=800000}}
\begin{tikzpicture}[femm/problem={xmin=0,xmax=42,ymin=-22,ymax=22,
    mesh size=1.5,depth=20},
    femm/pair={first=lower,second=upper,type=antiperiodic,rotation=60},
    femm/cache={file=\jobname-sector,mode=auto}]
  \draw[femm/domain] (-30:10)--(-30:40)
    arc[start angle=-30,end angle=30,radius=40]--(30:10)
    arc[start angle=30,end angle=-30,radius=10]--cycle;
  \filldraw[fill=red!20,femm/region={material=recoil,
      magnetization=angular,magnetization order=4}]
    (-30:12)--(-30:20) arc[start angle=-30,end angle=30,radius=20]
    --(30:12) arc[start angle=30,end angle=-30,radius=12]--cycle;
  \draw[paired edge=lower] (-30:10)--(-30:40);
  \draw[paired edge=upper] (30:10)--(30:40);
  \pic[femm/lines=16,femm/arrow length=.8mm,femm/arrow width=.5mm] {femm field};
  \node[align=center,font=\small,fill=white,inner sep=2pt] at (28,0) {One sector\\$A_2=-A_1$};
\end{tikzpicture}
\qquad
% Reuse the solved model. These outlines are drawing paths only.
\begin{tikzpicture}[x=1mm,y=1mm]
  \foreach \sector in {0,60,...,300}{
    \begin{scope}[rotate=\sector]
      \draw[gray!35] (-30:10)--(-30:40)
        arc[start angle=-30,end angle=30,radius=40]--(30:10)
        arc[start angle=30,end angle=-30,radius=10]--cycle;
      \filldraw[fill=red!12,draw=gray!50]
        (-30:12)--(-30:20) arc[start angle=-30,end angle=30,radius=20]
        --(30:12) arc[start angle=30,end angle=-30,radius=12]--cycle;
    \end{scope}}
  \pic[femm/repeat={copies=6},femm/lines=16,
    femm/arrow length=.8mm,femm/arrow width=.5mm] {femm repeated field};
  \node[font=\small] at (0,-44) {Six reconstructed sectors};
\end{tikzpicture}
\end{document}
